\newpage
\chapter{矩阵}
\section{矩阵的概念}
略.
\section{矩阵的运算}
4.计算下列矩阵的$k$次幂.
\begin{align*}
	&\textup{（1）}
	\bm{A}=
	\begin{pmatrix}
		a&&\\
		&b&\\
		&&c
	\end{pmatrix}
	\qquad\qquad\qquad\qquad
	\textup{（2）}
	\bm{A}=\begin{pmatrix}
		\cos \theta&-\sin\theta\\
		\sin\theta&\cos\theta
	\end{pmatrix}\\
	&\textup{（3）}
	\bm{A}=\begin{pmatrix}
		a&1&0\\
		0&a&1\\
		0&0&a
	\end{pmatrix}
	\qquad\qquad\qquad\qquad
	\textup{（4）}
	\bm{A}=\begin{pmatrix}
		1&2&4\\
		2&4&8\\
		3&6&12
	\end{pmatrix}
\end{align*}
\begin{solution}
	\textup{（1）}
	\[
	\bm{A}^k=\begin{pmatrix}
		a^k&&\\
		&b^k&\\
		&&c^k
	\end{pmatrix}
	\]

	\textup{（2）}
	考虑数学归纳法,或者利用几何意义,显然有
	\[
	\bm{A}^k=\begin{pmatrix}
		\cos k\theta &-\sin k\theta\\
		\sin k\theta &\cos k\theta
	\end{pmatrix}
	\]

	\textup{（3）}考虑拆分\[
	\bm{A}=a\bm{I}_3+\bm{J}
	\]
	其中
	\[
	\bm{J}=\begin{pmatrix}
		0&1&0\\0&0&1\\0&0&0
	\end{pmatrix}
	\]
	利用二项式定理展开即可
	$\left(\bm{J}_3=\bm{O}\right)$.
	\[
	\bm{A}^k=\begin{pmatrix}
		a^k&C_k^1a^{k-1}&C_k^2a^{k-2}\\
		0&a^k&C_k^1a^{k-1}\\
		0&0&a^k
	\end{pmatrix}\left(k\geqslant 2\right)
	\]

	\textup{（4）}
	注意到明显的对称性,设$\bm{\alpha}=\left(1,2,4
	\right),\bm{\beta}=\left(1,2,3\right)$,
	于是$\bm{A}=\bm{\beta}^T\bm{\alpha},
	\bm{\alpha}^T\bm{\beta}=17$.故
	\begin{align*}
			\bm{A}^k&=\left(\bm{A}\bm{A}\cdots\bm{A}\right)
\\
		&=\left(\bm{\beta}^T\bm{\alpha}\bm{\beta}^T\bm{\alpha}\cdots\bm{\beta}^T\bm{\alpha}\right)
		\\
		&=17^{k-1}\bm{A}
	\end{align*}
\end{solution}

5.设$n$阶对称阵$\bm{A}=\left(a_{ij}\right)$,
$\bm{x}=\left(x_1,x_2,\cdots,x_n\right)^T\in \mathbb{R}^n$,
求$\bm{x}^T\bm{Ax}.$
\begin{solution}
因为$a_{ij}=a_{ji}$
\begin{align*}
	&\begin{pmatrix}
		x_1,x_2,\cdots,x_n
	\end{pmatrix}\begin{pmatrix}
		a_{11}&a_{12}&\cdots&a_{1n}\\
		a_{21}&a_{22}&\cdots&a_{2n}\\
		\vdots&\vdots&&\vdots\\
		a_{n1}&a_{n2}&\cdots&a_{nn}
	\end{pmatrix}\begin{pmatrix}
		x_1\\
		x_2\\
		\vdots\\
		x_n
	\end{pmatrix}\\
	=&\begin{pmatrix}
		\sum\limits_{i=1}^na_{i1}x_i,
		\sum\limits_{i=1}^{n}a_{i2}x_i,\cdots,
		\sum\limits_{i=1}^{n}a_{in}x_i
	\end{pmatrix}
	\begin{pmatrix}
		x_1\\x_2\\\vdots\\x_n
	\end{pmatrix}
	\\
	=&\sum_{i=1}^{n}a_{ii}x_i^2+2\sum_{1\leqslant i<j\leqslant n}
	a_{ij}x_ix_j
\end{align*}
\end{solution}

8.设实对称矩阵$\bm{A}$,且$\bm{A}^2=\bm{O}$,
求证：$\bm{A}=\bm{O}$.
\begin{Proof}
	因为$a_{ij}=a_{ji}$,而
\[
\bm{A}^2=\bm{AA}^T=\left(a_{ij}^2\right)=\bm{O}\Longrightarrow a_{ij} = 0
\qedhere
\]
\end{Proof}

11.设$n$阶复矩阵$\bm{A}$,求证：
$\left|\overline{\bm{A}}\right|=
\overline{\left|\bm{A}\right|}$.
\begin{Proof}
考虑行列式的组合定义即可.
\end{Proof}

\section{方阵的逆阵}
4.设非异阵$\bm{A}$,求证：$\left(\bm{A}^k\right)^{-1}
=\left(\bm{A}^{-1}\right)^k.$
\begin{Proof}
$
\left(\bm{AA}\cdots\bm{A}\right)^{-1}=\left(\bm{A}^{-1}\bm{A}^{-1}\cdots\bm{A}^{-1}\right)=\left(\bm{A}^{-1}\right)^k
.$
\end{Proof}

6.设对合阵$\bm{A}$,且$\bm{A}+\bm{I}_n$为非异阵,求证：
$\bm{A}=\bm{I}_n.$
\begin{Proof}
$\left(\bm{A}+\bm{I}_n\right)\left(\bm{A}-\bm{I}_n\right)=\bm{O}\Longrightarrow
\bm{A}=\bm{I}_n.$
\end{Proof}

7.设非零实矩阵$\bm{A}$且$\bm{A}^*=\bm{A}^T$,求证：$\bm{A}$是非异阵.
\begin{Proof}
	已知$a_{ij}=\bm{A}_{ij}$.
显然,$\exists a_{rs}\neq 0$,那么展开得
\begin{align*}
	\left|\bm{A}\right|&=a_{r1}\bm{A}_{r1}+a_{r2}\bm{A}_{r2}+\cdots+a_{rn}\bm{A}_{rn}
\\
&=\sum_{i=1}^{n}a_{ri}^2
	>0
	\qedhere
\end{align*}
\end{Proof}

8.设$\bm{A}$、$\bm{B}$、$\bm{A}+\bm{B}$均为非异阵,
求证：
$\bm{A}^{-1}+\bm{B}^{-1}$非异.
\begin{Proof}
注意到
\[
\bm{A}\left(\bm{A}^{-1}+\bm{B}^{-1}\right)\bm{B}=\bm{A}+\bm{B}
\]
于是
\[
\bm{B}\left(\bm{A}+\bm{B}\right)^{-1}\bm{A}\left(\bm{A}^{-1}+\bm{B}^{-1}\right)=\bm{I}
\qedhere
\]
\end{Proof}

9.设幂零阵$\bm{A}$,求证：$\bm{I}_n-\bm{A}$非异.
\begin{Proof}
注意到
\[
\left(\bm{I}_n-\bm{A}\right)\left(
	\bm{I}_n+\bm{A}+\bm{A}^2+\cdots+\bm{A}^{m-1}\right)=\bm{I}_n
\qedhere
\]
\end{Proof}

10.设幂等阵$\bm{A}$,证明：$\bm{A}+\bm{I}_n$非异.
\begin{Proof}
注意到
\[
\left(\bm{A}+\bm{I}_n\right)\left(\bm{A}-\bm{I}_n\right)=-2\bm{I}_n
\qedhere
\]
\end{Proof}
\section{矩阵的初等变换与初等矩阵}
略.
\section{矩阵乘积的行列式与初等变换法求逆阵}
5.设\[
\bm{S}=\begin{pmatrix}
	s_0&s_1&s_2&\cdots&s_{n-1}\\
	s_1&s_2&s_3&\cdots&s_n\\
	s_2&s_3&s_4&\cdots&s_{n+1}\\
	\vdots&\vdots&\vdots&&\vdots\\
	s_{n-1}&s_n&s_{n+1}&\cdots&s_{2n-2}
\end{pmatrix}
\]
其中$s_k=x_1^k+x_2^k+\cdots+x_n^k.$求$\left|\bm{S}\right|$并证明
$\left|\bm{S}\right|\geqslant 0\left(\forall x_i\in\mathbb{R},i=1,2,\cdots,n\right)$.
\begin{Proof}
注意到
\begin{align*}
	\left|\bm{S}\right|&=\begin{vmatrix}
		\begin{pmatrix}
			1&1&\cdots&1&1\\
			x_1&x_2&\cdots&x_{n-1}&x_n\\
			\vdots&\vdots&&\vdots&\vdots\\
			x_1^{n-1}&x_2^{n-1}&\cdots&x_{n-1}^{n-1}&x_n^{n-1}
		\end{pmatrix}
		\begin{pmatrix}
			1&x_1&\cdots&x_1^{n-1}\\
			1&x_2&\cdots&x_2^{n-1}\\
			\vdots&\vdots&&\vdots\\
			1&x_n&\cdots&x_n^{n-1}
		\end{pmatrix}
	\end{vmatrix}\\
	&=\begin{vmatrix}
		1&1&\cdots&1&1\\
			x_1&x_2&\cdots&x_{n-1}&x_n\\
			\vdots&\vdots&&\vdots&\vdots\\
			x_1^{n-1}&x_2^{n-1}&\cdots&x_{n-1}^{n-1}&x_n^{n-1}
	\end{vmatrix}^2\\
	&=\prod_{1\leqslant i<j\leqslant n}
	\left(x_j-x_i\right)^2\geqslant 0\left(\forall x_k\in\mathbb{R},k=1,2,\cdots,
	n\right)\qedhere
\end{align*}
\end{Proof}

6.利用行列式乘法证明循环矩阵的行列式的值等于
$f\left(\varepsilon_1\right)f\left(\varepsilon_2\right)\cdots f\left(\varepsilon_n\right)$,其中
$f\left(x\right)=a_1+a_2x+a_3x^2+
\cdots+a_nx^{n-1},\varepsilon_i\left(i=1,2,\cdots,n\right)$为$1$的全体$n$次方根：
\[
\begin{vmatrix}
	a_1&a_2&a_3&\cdots&a_n\\
    a_n&a_1&a_2&\cdots&a_{n-1}\\
    a_{n-1}&a_n&a_1&\cdots&a_{n-2}\\
    \vdots&\vdots&\vdots&&\vdots\\
    a_2&a_3&a_4&\cdots&a_{1}
\end{vmatrix}
\]
\begin{Proof}
注意到\begin{align*}
	&\begin{pmatrix}
	a_1&a_2&a_3&\cdots&a_n\\
    a_n&a_1&a_2&\cdots&a_{n-1}\\
    a_{n-1}&a_n&a_1&\cdots&a_{n-2}\\
    \vdots&\vdots&\vdots&&\vdots\\
    a_2&a_3&a_4&\cdots&a_{1}
\end{pmatrix}\begin{pmatrix}
	1&1&\cdots&1&1\\
	\varepsilon_1&\varepsilon_2&\cdots&\varepsilon_{n-1}&\varepsilon_n\\
	\vdots&\vdots&&\vdots&\vdots\\
	\varepsilon_1^{n-1}&\varepsilon_2^{n-1}&\cdots&\varepsilon_{n-1}^{n-1}&\varepsilon_n^{n-1}
\end{pmatrix}\\
=&\begin{pmatrix}
	f\left(\varepsilon_1\right)&f\left(\varepsilon_2\right)&\cdots&f\left(\varepsilon_{n-1}\right)&f\left(\varepsilon_n\right)\\
	\varepsilon_1 f\left(\varepsilon_1\right)&\varepsilon_2 f\left(\varepsilon_2\right)&\cdots&\varepsilon_{n-1} f\left(\varepsilon_{n-1}\right)&\varepsilon_n f\left(\varepsilon_n\right)\\
	\vdots&\vdots&&\vdots&\vdots\\
	\varepsilon_1^{n-1} f\left(\varepsilon_1\right)&\varepsilon_2^{n-1} f\left(\varepsilon_2\right)&\cdots&\varepsilon_{n-1}^{n-1} f\left(\varepsilon_{n-1}\right)&\varepsilon_n ^{n-1} f\left(\varepsilon_n\right)
\end{pmatrix}\\
=&\begin{pmatrix}
		1&1&\cdots&1&1\\
	\varepsilon_1&\varepsilon_2&\cdots&\varepsilon_{n-1}&\varepsilon_n\\
	\vdots&\vdots&&\vdots&\vdots\\
	\varepsilon_1^{n-1}&\varepsilon_2^{n-1}&\cdots&\varepsilon_{n-1}^{n-1}&\varepsilon_n^{n-1}
	\end{pmatrix}\begin{pmatrix}
		f\left(\varepsilon_1\right)&&&\\
		&f\left(\varepsilon_2\right)&&\\
		&&\ddots&\\
		&&& f\left(\varepsilon_n\right)
	\end{pmatrix}
\end{align*}
于是
\begin{align*}
	&\begin{vmatrix}
		a_1&a_2&a_3&\cdots&a_n\\
    a_n&a_1&a_2&\cdots&a_{n-1}\\
    a_{n-1}&a_n&a_1&\cdots&a_{n-2}\\
    \vdots&\vdots&\vdots&&\vdots\\
    a_2&a_3&a_4&\cdots&a_{1}
	\end{vmatrix}
	\begin{vmatrix}
		1&1&\cdots&1&1\\
	\varepsilon_1&\varepsilon_2&\cdots&\varepsilon_{n-1}&\varepsilon_n\\
	\vdots&\vdots&&\vdots&\vdots\\
	\varepsilon_1^{n-1}&\varepsilon_2^{n-1}&\cdots&\varepsilon_{n-1}^{n-1}&\varepsilon_n^{n-1}
	\end{vmatrix}\\
	=&f\left(\varepsilon_1\right)f\left(\varepsilon_2\right)\cdots f\left(\varepsilon_n\right)
	\begin{vmatrix}
		1&1&\cdots&1&1\\
	\varepsilon_1&\varepsilon_2&\cdots&\varepsilon_{n-1}&\varepsilon_n\\
	\vdots&\vdots&&\vdots&\vdots\\
	\varepsilon_1^{n-1}&\varepsilon_2^{n-1}&\cdots&\varepsilon_{n-1}^{n-1}&\varepsilon_n^{n-1}
	\end{vmatrix}
\end{align*}
其中
\[
\begin{vmatrix}
		1&1&\cdots&1&1\\
	\varepsilon_1&\varepsilon_2&\cdots&\varepsilon_{n-1}&\varepsilon_n\\
	\vdots&\vdots&&\vdots&\vdots\\
	\varepsilon_1^{n-1}&\varepsilon_2^{n-1}&\cdots&\varepsilon_{n-1}^{n-1}&\varepsilon_n^{n-1}
	\end{vmatrix}=\prod_{1\leqslant i<j\leqslant n}\left(\varepsilon_j-\varepsilon_i\right)\neq 
	0\qedhere
\]
\end{Proof}

7.计算下面行列式的值：
\[
\left|\bm{A}\right|=\begin{vmatrix}
	1&2&3&\cdots&n-1&n\\
	n&1&2&\cdots&n-2&n-1\\
	n-1&n&1&\cdots&n-3&n-2\\
	\vdots&\vdots&\vdots&&\vdots&\vdots\\
	3&4&5&\cdots&1&2\\
	2&3&4&\cdots&n&1
\end{vmatrix}
\]
\begin{solution}
因为\begin{align*}
	f\left(x\right)&=1+2x+3x^2+\cdots+nx^{n-1}\\
	&=\left(x+x^2+\cdots+x^n\right)'
\\
&=
\left(\frac{x-x^{n+1}}{1-x}\right)'\\
&=\frac{1-\left(n+1\right)x^n+nx^{n+1}}{\left(1-x\right)^2}\left(x\neq 1
\right)
\end{align*}
于是\[
f\left(\varepsilon_i\right)=\begin{cases*}
\cfrac{n\left(n+1\right)}{2},i=n\\
\cfrac{n}{\varepsilon_i-1},1\leqslant i\leqslant n-1
\end{cases*}
\]
注意到\begin{align*}
	\prod_{i=1}^{n-1}\left(\varepsilon_i-1\right)&=f\left(z\right)=\frac{\left(\varepsilon_1-z\right)\left(\varepsilon_2-z\right)\cdots\left(\varepsilon_n-z\right)}{\varepsilon_n-z}\Bigg|_{z=1}\\
	&=\left(-1\right)^{n-1}\frac{z^n-1}{z-1}\Bigg|_{z=1}\\
	&=\left(-1\right)^{n-1}\left(1+z+z^2+\cdots+z^{n-1}\right)\Bigg|_{z=1}\\
	&=\left(-1\right)^{n-1}n
\end{align*}
于是\begin{align*}
	\left|\bm{A}\right|&=\frac{n\left(n+1\right)}{2}\frac{n^{n-1}}{\left(-1\right)^{n-1}n}\\
	&=\left(-1\right)^{n-1}\frac{\left(n+1\right)n^{n-1}}{2}
\end{align*}
\end{solution}
\section{分块矩阵}
3.设
\[
\bm{A}=\begin{pmatrix}
	0&1&0&\cdots&0&0\\
	0&0&1&\cdots&0&0\\
	\vdots&\vdots&\vdots&&\vdots&\vdots\\
	0&0&0&\cdots&0&1\\
	1&0&0&\cdots&0&0
\end{pmatrix}
\]
求证：
\[
\bm{A}^k=\begin{pmatrix}
	\bm{O}&\bm{I}_{n-k}\\
	\bm{I}_k&\bm{O}
\end{pmatrix}\left(k=1,2,\cdots,n\right)
\]
\begin{Proof}
注意到\begin{align*}
	\bm{A}^2&=\bm{A}\left(\bm{e}_n,\bm{e}_1,\bm{e}_2,\cdots,\bm{e}_{n-1}\right)\\
	&=\left(\bm{A}\bm{e}_n,\bm{A}\bm{e}_1,\bm{A}\bm{e}_2,\cdots,\bm{A}\bm{e}_{n-1}\right)\\
	&=\left(\bm{e}_{n-1},\bm{e}_n,\bm{e}_1,\cdots,\bm{e}_{n-2}\right)\qedhere
\end{align*}
\end{Proof}

8.设$n$阶方阵$\bm{A}$、$\bm{B}$,求证：
\[
\begin{vmatrix}
	\bm{A}&\bm{B}\\\bm{B}&\bm{A}
\end{vmatrix}=\left|\bm{A}+\bm{B}\right|\left|\bm{A}-\bm{B}\right|
\]
\begin{Proof}
\begin{align*}
	LHS&=\begin{vmatrix}
		\bm{A}+\bm{B}&\bm{A}+\bm{B}\\\bm{B}&\bm{A}
	\end{vmatrix}\\
	&=\begin{vmatrix}
		\bm{A}+\bm{B}&\bm{O}\\\bm{B}&\bm{A}-\bm{B}
	\end{vmatrix}\\
	&=\left|\bm{A}+\bm{B}\right|\left|\bm{A}-\bm{B}\right|\\
	&=RHS\qedhere
\end{align*}
\end{Proof}

9.设$n$阶复方阵$\bm{A}$、$\bm{B}$,求证：
\[
\begin{vmatrix}
	\bm{A}&-\bm{B}\\\bm{B}&\bm{A}
\end{vmatrix}=\left|\bm{A}+\mathrm{i}\bm{B}\right|\left|\bm{A}-\mathrm{i}\bm{B}\right|
\]
\begin{Proof}
\begin{align*}
	LHS&=\begin{vmatrix}
		\bm{A}+\mathrm{i}\bm{B}&-\left(\bm{B}-\mathrm{i}\bm{A}\right)\\
		\bm{B}&\bm{A}
	\end{vmatrix}\\
	&=\begin{vmatrix}
		\bm{A}+\mathrm{i}\bm{B}&\bm{O}\\
		\bm{B}&\bm{A}-\mathrm{i}\bm{B}
	\end{vmatrix}\\
	&=\left|\bm{A}+\mathrm{i}\bm{B}\right|\left|\bm{A}-\mathrm{i}\bm{B}\right|\\
	&=RHS\qedhere
\end{align*}
\end{Proof}

10.设$n$阶方阵$\bm{A}$、$\bm{B}$满足$\bm{AB}=\bm{BA}$,求证：
\[
\begin{vmatrix}
	\bm{A}&-\bm{B}\\\bm{B}&\bm{A}
\end{vmatrix}=\left|\bm{A}^2+\bm{B}^2\right|
\]
\begin{Proof}
\begin{align*}
	LHS&=\left|\bm{A}+\mathrm{i}\bm{B}\right|\left|\bm{A}-\mathrm{i}\bm{B}\right|
	\\
	&=\left|\left(\bm{A}+\mathrm{i}\bm{B}\right)\left(\bm{A}-\mathrm{i}\bm{B}\right)\right|\\
	&=\left|\bm{A}^2+\bm{B}^2-\mathrm{i}\left(\bm{AB}-\bm{BA}\right)\right|\\
	&=\left|\bm{A}^2+\bm{B}^2\right|\\
	&=RHS\qedhere
\end{align*}
\end{Proof}

12.求下列矩阵$\bm{A}$的行列式的值：
\[
(1)\bm{A}=\begin{pmatrix}
	0&2&3&\cdots&n\\
	1&0&3&\cdots&n\\
	1&2&0&\cdots&n\\
	\vdots&\vdots&\vdots&&\vdots\\
	1&2&3&\cdots&0
\end{pmatrix}\qquad\qquad\qquad
(2)\bm{A}=\begin{pmatrix}
	1+a_1^2&a_1a_2&\cdots&a_1a_n\\
	a_2a_1&1+a_2^2&\cdots&a_2a_n\\
	\vdots&\vdots&&\vdots\\
	a_na_1&a_na_2&\cdots&1+a_n^2
\end{pmatrix}
\]
\begin{solution}
$(1)$注意到\begin{align*}
\bm{A}&=\begin{pmatrix}
	-1&0&\cdots&0\\
	0&-2&\cdots&0\\
	\vdots&\vdots&&\vdots\\
	0&0&\cdots&-n
\end{pmatrix}+\begin{pmatrix}
	1\\1\\\vdots\\1
\end{pmatrix}\left(1,2,\cdots,n\right)
\end{align*}
因此由降阶公式得
\begin{align*}
	\left|\bm{A}\right|&=\begin{vmatrix}
		-1&0&\cdots&0\\
	0&-2&\cdots&0\\
	\vdots&\vdots&&\vdots\\
	0&0&\cdots&-n
	\end{vmatrix}\left|
		1+\left(1,2,\cdots,n\right)\begin{pmatrix}
			-1&0&\cdots&0\\
	0&-2&\cdots&0\\
	\vdots&\vdots&&\vdots\\
	0&0&\cdots&-n
		\end{pmatrix}^{-1}\begin{pmatrix}
			1\\1\\\vdots\\1
		\end{pmatrix}
	\right|\\
	&=\left(-1\right)^nn!\left(1-n\right)
\end{align*}

$(2)$注意到\begin{align*}
	\bm{A}=\bm{I}_n+\begin{pmatrix}
		a_1\\a_2\\\vdots\\a_n
	\end{pmatrix}\left(a_1,a_2,\cdots,a_n\right)
\end{align*}
故由降阶公式得
\begin{align*}
	\left|\bm{A}\right|&=
	\left|1+\left(a_1,a_2,\cdots,a_n\right)\begin{pmatrix}
		a_1\\a_2\\\vdots\\a_n
	\end{pmatrix}\right|\\
	&=1+\sum_{i=1}^{n}a_i^2
\end{align*}
\end{solution}
\section{Cauchy-Binet公式}
1.设$n\geqslant 3$,求证：
\[
\left|\bm{A}\right|=\begin{vmatrix}
    1+x_1y_1&1+x_1y_2&\cdots&1+x_1y_n\\
    1+x_2y_1&1+x_2y_2&\cdots&1+x_2y_n\\
    \vdots&\vdots&&\vdots\\
    1+x_ny_1&1+x_ny_2&\cdots&1+x_ny_n
\end{vmatrix}=0
\]
\begin{Proof}
注意到
\[
\bm{A}=\begin{bmatrix}
    1&x_1&0&\cdots&0\\
    1&x_2&0&\cdots&0\\
    \vdots&\vdots&\vdots&&\vdots\\
    1&x_n&0&\cdots&0 
\end{bmatrix}
\begin{bmatrix}
    1&1&\cdots&1\\
    y_1&y_2&\cdots&y_n\\
    0&0&\cdots&0\\
    \vdots&\vdots&&\vdots\\
    0&0&\cdots&0
\end{bmatrix}\qedhere
\]
\end{Proof}
2.设$n$阶实方阵$\bm{A}$满足$\bm{AA}'=\bm{I}_n$,求证：$\forall 1\leqslant i_1<i_2<\cdots<i_r\leqslant n$
\[
\sum_{
    1\leqslant j_1<j_2<\cdots<j_r\leqslant n
}\bm{A}\begin{pmatrix}
    i_1&i_2&\cdots&i_r\\
    j_1&j_2&\cdots&j_r
\end{pmatrix}^2=1
\]
\begin{Proof}
对$\bm{AA}'=\bm{I}_n$两边同时求行列式得
\begin{align*}
     LHS &=\sum_{
        1\leqslant j_1<j_2<\cdots<j_r\leqslant n
    }\bm{A}\begin{pmatrix}
        i_1&i_2&\cdots&i_r\\
        j_1&j_2&\cdots&j_r
    \end{pmatrix}
    \bm{A}'\begin{pmatrix}
        j_1&j_2&\cdots&j_r\\
        i_1&i_2&\cdots&i_r
    \end{pmatrix}\\
    &=\sum_{
        1\leqslant j_1<j_2<\cdots<j_r\leqslant n
    }\bm{A}\begin{pmatrix}
        i_1&i_2&\cdots&i_r\\
        j_1&j_2&\cdots&j_r
    \end{pmatrix}^2\\
    &=1=RHS\qedhere
\end{align*}
\end{Proof}
3.设$n$阶方阵$\bm{A},\bm{B}$,$r$是不大于$n$的正整数.求证：$\bm{AB}$和$\bm{BA}$的所有$r$阶主子式之和相等.
\begin{Proof}
\begin{align*}
    &\sum_{
        1\leqslant i_1<i_2<\cdots<i_r\leqslant n
    }\bm{AB}\begin{pmatrix}
        i_1&i_2&\cdots&i_r\\
        i_1&i_2&\cdots&i_r
    \end{pmatrix}\\
    &=\sum_{
        1\leqslant i_1<i_2<\cdots<i_r\leqslant n
    }\sum_{
        1\leqslant j_1<j_2<\cdots<j_r\leqslant n
    }\bm{A}\begin{pmatrix}
        i_1&i_2&\cdots&i_r\\
        j_1&j_2&\cdots&j_r
    \end{pmatrix}
    \bm{B}\begin{pmatrix}
        j_1&j_2&\cdots&j_r\\
        i_1&i_2&\cdots&i_r
    \end{pmatrix}\\
    &=\sum_{
        1\leqslant j_1<j_2<\cdots<j_r\leqslant n
    }
    \sum_{
        1\leqslant i_1<i_2<\cdots<i_r\leqslant n
    }
    \bm{B}\begin{pmatrix}
        j_1&j_2&\cdots&j_r\\
        i_1&i_2&\cdots&i_r
    \end{pmatrix}
    \bm{A}\begin{pmatrix}
        i_1&i_2&\cdots&i_r\\
        j_1&j_2&\cdots&j_r
    \end{pmatrix}\\
    &=\sum_{
        1\leqslant j_1<j_2<\cdots<j_r\leqslant n
    }\bm{BA}\begin{pmatrix}
        j_1&j_2&\cdots&j_r\\
        j_1&j_2&\cdots&j_r
    \end{pmatrix}\qedhere
\end{align*}
\end{Proof}
4.设$m\times n$阶实矩阵$\bm{A},\bm{B}$,求证：
\[
\left|
    \bm{AA}'
\right|\left|
    \bm{BB}'
\right|\geqslant\left|
    \bm{AB}'
\right|^2
\]
\begin{Proof}
一方面,$m > n$时,$\left|
    \bm{AA}'
\right|=\left|
    \bm{BB}'
\right|=0=\left|
    \bm{AB}'
\right|$,结论显然成立.

另一方面,$m\leqslant n$时
\begin{align*}
    &LHS\\
    =&\sum_{
        1\leqslant i_1<i_2<\cdots<i_m\leqslant n
    }\bm{A}\begin{pmatrix}
        1&2&\cdots&m\\
        i_1&i_2&\cdots&i_m
    \end{pmatrix}^2\sum_{
        1\leqslant j_1<j_2<\cdots<j_m\leqslant n
    }\bm{B}\begin{pmatrix}
        1&2&\cdots&m\\
        j_1&j_2&\cdots&j_m
    \end{pmatrix}^2\\
    &\geqslant\left(
        \sum_{
            1\leqslant i_1<i_2<\cdots<i_m\leqslant n
        }\bm{A}\begin{pmatrix}
            1&2&\cdots&m\\
            i_1&i_2&\cdots&i_m
        \end{pmatrix}\bm{B}\begin{pmatrix}
            1&2&\cdots&m\\
            i_1&i_2&\cdots&i_m
        \end{pmatrix}
    \right)      ^2=RHS\qedhere               
\end{align*}
\end{Proof}
5.设$m\times n$阶复矩阵$\bm{A}$,求证：
$\bm{A}\overline{\bm{A}}'$的任一主子式非负.
\begin{Proof}
一方面,$r>n$时$\bm{A}\overline{\bm{A}}'$的$r$阶主子式为$0$,结论显然成立.

另一方面,$r\leqslant n$时
\begin{align*}
    &\bm{A}\overline{\bm{A}}'\begin{pmatrix}
        i_1&i_2&\cdots&i_r\\
        i_1&i_2&\cdots&i_r
    \end{pmatrix}\\
    =&
    \sum_{
        1\leqslant j_1<j_2<\cdots<j_r\leqslant n
    }\bm{A}\begin{pmatrix}
        i_1&i_2&\cdots&i_r\\
        j_1&j_2&\cdots&j_r
    \end{pmatrix}
    \overline{\bm{A}}'\begin{pmatrix}
        j_1&j_2&\cdots&j_r\\
        i_1&i_2&\cdots&i_r
    \end{pmatrix}\\
    =&\sum_{
        1\leqslant j_1<j_2<\cdots<j_r\leqslant n
    }\bm{A}\begin{pmatrix}
        i_1&i_2&\cdots&i_r\\
        j_1&j_2&\cdots&j_r
    \end{pmatrix}
    \overline{\bm{A}}\begin{pmatrix}
        i_1&i_2&\cdots&i_r\\
        j_1&j_2&\cdots&j_r
    \end{pmatrix}\\
    =&\sum _{
        1\leqslant j_1<j_2<\cdots<j_r\leqslant n
    }\bm{A}\begin{pmatrix}
        i_1&i_2&\cdots&i_r\\
        j_1&j_2&\cdots&j_r
    \end{pmatrix}\overline{
          \bm{A}\begin{pmatrix}
                i_1&i_2&\cdots&i_r\\
                j_1&j_2&\cdots&j_r
            \end{pmatrix}   
    }\\
    =&\sum_{
        1\leqslant j_1<j_2<\cdots<j_r\leqslant n
    }\left|\bm{A}\begin{pmatrix}
        i_1&i_2&\cdots&i_r\\
        j_1&j_2&\cdots&j_r
    \end{pmatrix}\right|^2\geqslant 0\qedhere
\end{align*}
\end{Proof}
6.证明复数形式的Cauchy-Schwarz不等式：
\[
\left|\sum_{i=1}^{n}a_i\overline{b}_i\right|^2\leqslant\sum_{i=1}^{n}\left|a_i\right|^2\sum_{i=1}^{n}\left|b_i\right|^2
\]
其中$a_i,b_i\in\mathbb{C},n\geqslant 2.$
\begin{Proof}
考虑如下行列式
\[
\begin{vmatrix}
    \sum\limits_{i=1}^{n}\left|a_i\right|^2&
    \sum\limits_{i=1}^{n}a_i\overline{b}_i\\
    \sum\limits_{i=1}^{n}\overline{a}_ib_i&
    \sum\limits_{i=1}^{n}\left|b_i\right|^2
\end{vmatrix}=
\left|
    \begin{pmatrix}
        a_1&a_2&\cdots&a_n\\
        b_1&b_2&\cdots&b_n 
    \end{pmatrix}
    \begin{pmatrix}
        \overline{a}_1&\overline{b}_1\\
        \overline{a}_2&\overline{b}_2\\
        \vdots&\vdots\\
        \overline{a}_n&\overline{b}_n
    \end{pmatrix}
\right|
\]
由Cauchy-Binet公式得
\begin{align*}
    &\left|
        \begin{pmatrix}
            a_1&a_2&\cdots&a_n\\
            b_1&b_2&\cdots&b_n 
        \end{pmatrix}
        \begin{pmatrix}
            \overline{a}_1&\overline{b}_1\\
            \overline{a}_2&\overline{b}_2\\
            \vdots&\vdots\\
            \overline{a}_n&\overline{b}_n
        \end{pmatrix}
    \right|\\
    &=\sum_{1\leqslant i<j\leqslant n}
    \begin{vmatrix}
        a_i&a_j\\
        b_i&b_j
    \end{vmatrix}
    \begin{vmatrix}
        \overline{a}_i&\overline{a}_j\\
        \overline{b}_i&\overline{b}_j
    \end{vmatrix}\\
    &=\sum_{1\leqslant i<j\leqslant n}
    \left(
        a_ib_j-a_jb_i
    \right)\left(
        \overline{a}_i\overline{b}_j-\overline{a}_j\overline{b}_i
    \right)\\
    &=\sum_{1\leqslant i<j\leqslant n}
    \left|
        a_ib_j-a_jb_i 
    \right|^2\geqslant 0\qedhere
\end{align*}
\end{Proof}
7.证明推广到复数形式的第4题.
\begin{Proof}
即设$m\times n$阶复矩阵$\bm{A},\bm{B}$,求证：
\[
\left|
    \bm{A}\overline{\bm{A}}'
\right|\left|
    \bm{B}\overline{\bm{B}}'
\right|\geqslant
\left|
    \bm{A}\overline{\bm{B}}'
\right|^2
\]
一方面,$m>n$时$\left|
    \bm{A}\overline{\bm{A}}'
\right|=\left|
    \bm{B}\overline{\bm{B}}'
\right|=0=\left|
    \bm{A}\overline{\bm{B}}'
\right|$,结论显然成立.

另一方面,$m\leqslant n$时
\begin{align*}
    &LHS\\
    =&\sum_{
        1\leqslant i_1<i_2<\cdots<i_m\leqslant n
    }\bm{A}\begin{pmatrix}
        1&2&\cdots&m\\
        i_1&i_2&\cdots&i_m
    \end{pmatrix}\overline{\bm{A}}\begin{pmatrix}
        1&2&\cdots&m\\
        i_1&i_2&\cdots&i_m
    \end{pmatrix}\\
    \cdot &\sum_{
        1\leqslant j_1<j_2<\cdots<j_m\leqslant n
    }\bm{B}\begin{pmatrix}
        1&2&\cdots&m\\
        j_1&j_2&\cdots&j_m
    \end{pmatrix}\overline{\bm{B}}\begin{pmatrix}
        1&2&\cdots&m\\
        j_1&j_2&\cdots&j_m
    \end{pmatrix}\\
    =&\sum _{
        1\leqslant i_1<i_2<\cdots<i_m\leqslant n
    }\left|
        \bm{A}\begin{pmatrix}
            1&2&\cdots&m\\
            i_1&i_2&\cdots&i_m 
        \end{pmatrix}
    \right|^2\sum _{
        1\leqslant i_1<i_2<\cdots<i_m\leqslant n
    }\left|
        \overline{\bm{B}}\begin{pmatrix}
            1&2&\cdots&m\\
            i_1&i_2&\cdots&i_m 
        \end{pmatrix}
    \right|^2\\
    \geqslant&\left(
        \sum_{
            1\leqslant i_1<i_2<\cdots<i_m\leqslant n
        }\left|
            \bm{A}\begin{pmatrix}
                1&2&\cdots&m\\
                i_1&i_2&\cdots&i_m 
            \end{pmatrix}
            \overline{\bm{B}}\begin{pmatrix}
                1&2&\cdots&m\\
                i_1&i_2&\cdots&i_m 
            \end{pmatrix}  
        \right|
    \right)^2\\
    =&RHS\qedhere
\end{align*}
\end{Proof}